% ============================================================
%  La ecuación de la recta — carrusel 1:1
%  Estilo: mathwizards-carrusel-moderno
%  Compilar:  latexmk -xelatex ecuacion_recta.tex
% ============================================================
\documentclass[12pt]{article}
\usepackage{mathwizards-carrusel-moderno}
\mwcartotal{10}

% ── Mini-gráficas reutilizables ──────────────────────────────
% Ejes cartesianos compactos (x,y en cm; escala 0.72 cm/unidad)
\newcommand{\ejesbase}{%
  \draw[mwink!25, line width=0.5pt, -{Stealth[length=1.4mm]}]
    (0,0) -- (4.55,0);
  \draw[mwink!25, line width=0.5pt, -{Stealth[length=1.4mm]}]
    (0,0) -- (0,4.35);
  \node[anchor=north east, font=\fontsize{6}{6}\selectfont, text=mwink!45]
    at (4.55,0) {$x$};
  \node[anchor=north east, font=\fontsize{6}{6}\selectfont, text=mwink!45]
    at (0,4.35) {$y$};
}

\begin{document}

% ── 1. Portada ───────────────────────────────────────────────
\mwmodportada
  {Álgebra · Funciones lineales}
  {La ecuación de la recta}
  {Qué significa, cómo se escribe y cómo se usa para resolver
   problemas del mundo real.}

% ── 2. Anatomía de la ecuación ───────────────────────────────
\begin{mwmodlamina}[mwprimario]{La fórmula}{Una recta, tres letras}
\begin{mwcajaformula}
  \scalebox{1.12}{$\displaystyle
    y = \underbrace{m}_{\text{\fontsize{7}{8}\selectfont pendiente}}\;x
      + \underbrace{b}_{\text{\fontsize{7}{8}\selectfont intercepto}}$}
\end{mwcajaformula}

\begin{mwcaja}
  \mwctitulo{Qué hace cada letra}
  \noindent
  \begin{minipage}[t]{4.1cm}
    {\nxheav\fontsize{11}{11}\selectfont\color{mwprimario}$m$}\quad
    {\mwsmall mide la inclinación: cuánto cambia $y$ por cada
    unidad de $x$.}
  \end{minipage}\hfill
  \begin{minipage}[t]{4.1cm}
    {\nxheav\fontsize{11}{11}\selectfont\color{mwprimario}$b$}\quad
    {\mwsmall es el punto de partida: el valor de $y$ cuando
    $x=0$.}
  \end{minipage}
\end{mwcaja}

\begin{mwcajagris}
  {\mwsmall $x$ e $y$ son las coordenadas de \emph{cualquier}
  punto de la recta: si un par las cumple, el punto está en la recta.}
\end{mwcajagris}
\end{mwmodlamina}

% ── 3. La pendiente ──────────────────────────────────────────
\begin{mwmodlamina}[mwmagenta]{Pendiente}{La pendiente: cuánto sube}
\noindent
\begin{minipage}[t]{4.05cm}
  \vspace*{0pt}
  \begin{tikzpicture}[x=0.78cm,y=0.78cm]
    \ejesbase
    \draw[mwmagenta, line width=1.6pt] (0.35,0.55) -- (3.85,3.6);
    \draw[mwmagenta!60, dashed, line width=0.7pt]
      (1.0,1.15) -- (3.2,1.15) -- (3.2,3.05);
    \node[font=\fontsize{6.4}{6.4}\selectfont, text=mwmagenta]
      at (2.1,0.8) {$\Delta x$};
    \node[font=\fontsize{6.4}{6.4}\selectfont, text=mwmagenta, rotate=90]
      at (3.4,2.1) {$\Delta y$};
    \fill[mwmagenta] (1.0,1.15) circle (1.4pt);
    \fill[mwmagenta] (3.2,3.05) circle (1.4pt);
  \end{tikzpicture}
  \par\vspace{2pt}
  {\mwsmall\color{mwink!55}
   $\Delta y$ sube \;·\; $\Delta x$ avanza}
\end{minipage}\hfill
\begin{minipage}[t]{4.7cm}
  \vspace*{0pt}
  \begin{mwcaja}
    \mwctitulo{Cómo se calcula}
    \centering
    $m=\dfrac{\Delta y}{\Delta x}
      =\dfrac{y_2-y_1}{x_2-x_1}$
  \end{mwcaja}
  \begin{mwcaja}
    \mwctitulo{Lectura rápida}
    {\mwsmall
    $m>0$: la recta sube \quad\newline
    $m<0$: la recta baja \quad\newline
    $m=0$: la recta es horizontal}
  \end{mwcaja}
\end{minipage}

\vspace{2pt}
\begin{mwcajadato}
  {\mwsmall La pendiente es una \textbf{razón de cambio}: por cada
  unidad que avanzas, la recta sube o baja $m$ unidades.}
\end{mwcajadato}
\end{mwmodlamina}

% ── 4. El intercepto ─────────────────────────────────────────
\begin{mwmodlamina}[mwnaranja]{Intercepto}{El punto de partida: $b$}
\noindent
\begin{minipage}[t]{4.7cm}
  \vspace*{0pt}
  \begin{mwcaja}
    \mwctitulo{Qué es}
    {\mwsmall Es el valor de $y$ cuando $x=0$. La recta corta al
    eje $y$ en el punto $(0,b)$.}
  \end{mwcaja}
  \begin{mwcaja}
    \mwctitulo{Ejemplo}
    {\mwsmall En $y=3x+2$, la recta corta al eje $y$ en
    $(0,2)$.}
  \end{mwcaja}
  \begin{mwcajaalerta}
    {\mwsmall $b$ \textbf{no} es la raíz: la raíz es donde
    $y=0$, no donde $x=0$.}
  \end{mwcajaalerta}
\end{minipage}\hfill
\begin{minipage}[t]{4.05cm}
  \vspace*{0pt}
  \begin{mwcajaformula}
    \begin{tikzpicture}[x=0.64cm,y=0.64cm]
      \ejesbase
      \draw[mwnaranja, line width=1.6pt] (0,1.35) -- (3.9,3.9);
      \fill[mwnaranja] (0,1.35) circle (1.5pt);
      \draw[mwnaranja!55, dashed, line width=0.7pt] (0,1.35) -- (0.95,1.35);
      \node[anchor=west, font=\fontsize{6.4}{6.4}\selectfont, text=mwnaranja]
        at (0.08,1.05) {$(0,b)$};
    \end{tikzpicture}
  \end{mwcajaformula}
  \par\vspace{1pt}
  {\mwsmall\color{mwink!55}Corta al eje $y$ en $(0,b)$.}
\end{minipage}
\end{mwmodlamina}

% ── 5. Tres formas ───────────────────────────────────────────
\begin{mwmodlamina}[mwvioleta]{Formas}{Tres formas, una misma recta}
\begin{mwcaja}
  \mwctitulo{Elige según los datos que tengas}
  \mwcpill{mwvioleta}{Punto-pendiente}\quad
    $y-y_1=m(x-x_1)$\\[6pt]
  \mwcpill{mwprimario}{Explícita}\quad
    $y=mx+b$\\[6pt]
  \mwcpill{mwmagenta}{General}\quad
    $Ax+By+C=0$
\end{mwcaja}

\begin{mwcajapaso}
  {\mwsmall
  \textbf{Dos puntos} $\rightarrow$ calcula $m$ y usa
  punto-pendiente.\newline
  \textbf{Pendiente y corte} $\rightarrow$ escribe directo
  $y=mx+b$.\newline
  \textbf{Datos mezclados} $\rightarrow$ despeja $y$ y compara.}
\end{mwcajapaso}

\begin{mwcajagris}
  {\mwsmall Las tres describen \emph{la misma} recta: cambia la
  forma, no la gráfica.}
\end{mwcajagris}
\end{mwmodlamina}

% ── 6. Usos ──────────────────────────────────────────────────
\begin{mwmodlamina}[mwnaranja]{Aplicaciones}{¿Para qué sirve una recta?}
\begin{mwcaja}
  \mwcnum{mwnaranja}{1}\quad
  \textbf{Modelar costos.}\quad
  {\mwsmall $C(x)=(\text{tasa})\,x+(\text{fijo})$}\\[7pt]
  \mwcnum{mwnaranja}{2}\quad
  \textbf{Describir cambios.}\quad
  {\mwsmall la pendiente es la velocidad, la tasa o el ritmo.}\\[7pt]
  \mwcnum{mwnaranja}{3}\quad
  \textbf{Predecir.}\quad
  {\mwsmall estima valores que no mediste, dentro del modelo.}
\end{mwcaja}

\begin{mwcaja}
  \mwctitulo{La idea de fondo}
  {\mwsmall Si algo cambia \textbf{a tasa constante}, hay una
  recta detrás. Y si hay una recta, puedes calcular, comparar
  y predecir.}
\end{mwcaja}

\begin{mwcajaformula}
  \scalebox{0.95}{$\displaystyle
    \text{cambio total}=m\cdot(\text{cambio en }x)$}
\end{mwcajaformula}
\end{mwmodlamina}

% ── 7. Ejemplo: planteamiento ────────────────────────────────
\begin{mwmodlamina}[mwprimario]{Ejemplo real}{El taxi: armar el modelo}
\begin{mwcajapaso}
  {\mwsmall
  Un taxi cobra \textbf{\$1.500} de bajada de bandera y
  \textbf{\$600} por cada kilómetro recorrido.}
\end{mwcajapaso}

\begin{mwcaja}
  \mwctitulo{Traduce cada dato}
  {\mwsmall
  Variable $x$: kilómetros recorridos.\\[3pt]
  Costo fijo $b$: \$1.500 (se paga aunque $x=0$).\\[3pt]
  Tasa $m$: \$600 por km.}
\end{mwcaja}

\begin{mwcajaformula}
  \scalebox{1.1}{$\displaystyle C(x)=600x+1\,500$}
\end{mwcajaformula}

\begin{mwcajaalerta}
  {\mwsmall \textbf{Pregunta:} ¿cuánto cuesta un viaje de
  \textbf{8 km}?}
\end{mwcajaalerta}
\end{mwmodlamina}

% ── 8. Ejemplo: resolución ───────────────────────────────────
\begin{mwmodlamina}[mwmagenta]{Ejemplo real}{El taxi: paso a paso}
\begin{mwcajapaso}
  \mwctitulo{Sustituye $x=8$}
  $C(8)=600\,(8)+1\,500$\\[3pt]
  $C(8)=4\,800+1\,500$\\[3pt]
  $C(8)=6\,300$
\end{mwcajapaso}

\begin{mwcajaok}
  \mwctitulo{Respuesta}
  {\nxheav\fontsize{12}{14}\selectfont
  Un viaje de 8 km cuesta \textbf{\$6.300}.}
\end{mwcajaok}

\begin{mwcaja}
  \mwctitulo{Cómo leer el resultado}
  {\mwsmall La pendiente $m=600$ dice que \emph{cada kilómetro
  extra} suma \$600. El intercepto $b=1\,500$ es lo que pagas
  incluso antes de avanzar.}
\end{mwcaja}
\end{mwmodlamina}

% ── 9. Errores comunes ───────────────────────────────────────
\begin{mwmodlamina}[mwrojonaranja]{Ojo con esto}{3 errores que cuestan puntos}
\begin{mwcajaalerta}
  {\textbf{1. Invertir la pendiente.}\quad
  $m=\dfrac{\Delta x}{\Delta y}$ \;\mwcmal\\[2pt]
  {\mwsmall Correcto: } $m=\dfrac{\Delta y}{\Delta x}$
  \;\mwcbien}
\end{mwcajaalerta}

\begin{mwcajaalerta}
  {\textbf{2. Confundir $b$ con la raíz.}\quad
  \mwcmal\\[2pt]
  {\mwsmall $b$ es el corte en $y$; la raíz es el corte en $x$
  (donde $y=0$).}}
\end{mwcajaalerta}

\begin{mwcajaalerta}
  {\textbf{3. Restar en distinto orden.}\quad
  $\dfrac{y_2-y_1}{x_1-x_2}$ \;\mwcmal\\[2pt]
  {\mwsmall Arriba y abajo, \emph{el mismo} orden: }
  $\dfrac{y_2-y_1}{x_2-x_1}$ \;\mwcbien}
\end{mwcajaalerta}
\end{mwmodlamina}

% ── 10. Cierre ───────────────────────────────────────────────
\mwmodfin
  {Para llevar}
  {Una recta es la forma más simple\\
   de describir cómo cambia algo.}
  {Guárdalo y practica: elige dos datos\\
   de tu día y busca su recta.}

\end{document}
